% 03-copy-task.tex ;  The copy task, expanded with per-step norms and figure
%
% Scope: definition, experimental setup, per-step gradient norms,
% gradient decay curve figure, interpretation

\section{Bài toán sao chép}
\label{sec:ch02-copy}

\index{copy task|(}%

\subsection{Định nghĩa}

Bài toán sao chép kiểm tra trí nhớ theo cách trực tiếp nhất: mạng đọc một
chuỗi, im lặng qua một khoảng trễ, rồi viết lại chính xác chuỗi đó.

\begin{figure}[ht]
  \centering
  % ch02-copy-task-schema.tex — Copy task architecture
%
% Shows the three phases: input sequence → delimiter → blank/output.

\begin{tikzpicture}[
    >=Stealth,
    node distance=0.6cm and 0.0cm,
    every node/.style={font=\footnotesize},
    box/.style={draw, minimum width=0.75cm, minimum height=0.75cm},
    blank/.style={draw, minimum width=0.75cm, minimum height=0.75cm, fill=black!5},
    delim/.style={draw, minimum width=0.75cm, minimum height=0.75cm, fill=black!20},
    label/.style={font=\small\bfseries},
  ]

  % Phase 1: Input
  \node[label, above=0.3cm] at (1.5, 0) {Đọc vào (input)};
  \foreach \i/\v in {1/A, 2/B, 3/C} {
    \node[box] (in\i) at (\i*0.85, 0) {\v};
  }

  % Delimiter
  \node[label, above=0.3cm] at (4.3, 0) {Delimiter};
  \node[delim] (delim) at (4.3, 0) {$|$};

  % Phase 2: Blank
  \node[label, above=0.3cm] at (7.0, 0) {Khoảng trống};
  \foreach \i in {1,2,3} {
    \node[blank] (blank\i) at (5.0+\i*0.85, 0) {0};
  }

  % Output labels (below)
  \foreach \i/\v in {1/A, 2/B, 3/C} {
    \node[box, draw=black!40] (out\i) at (5.0+\i*0.85, -1.5) {\v};
  }

  % Arrows showing the RNN reads, then writes
  \draw[->, thick, black!60] (in1.south) to[bend right=20] (out1.north);
  \draw[->, thick, black!60] (in2.south) to[bend right=25] (out2.north);
  \draw[->, thick, black!60] (in3.south) to[bend right=30] (out3.north);

  % Phase labels
  \node[below=0.2cm of out2, font=\small\bfseries] {Viết ra (output)};

  % Arrow from input to output through RNN memory
  \draw[->, thick, black!40]
    (in3.east) -- (delim.west) node[midway, above=0.1cm, font=\tiny] {RNN nhớ};
  \draw[->, thick, black!40]
    (delim.east) -- (blank1.west) node[midway, above=0.1cm, font=\tiny] {nhớ tiếp};

  % Legend
  \node[below=3cm of blank2, font=\footnotesize\itshape, black!50, text width=6cm, align=center]
    {RNN đọc A, B, C; bắt đầu chép sau dấu $|$;\\
     phải nhớ toàn bộ chuỗi qua $T_{\text{mem}}$ bước im lặng.};

\end{tikzpicture}

  \caption{Sơ đồ bài toán sao chép với $T_{\text{mem}} = 3$. Mạng đọc A, B, C;
    sau delimiter $|$, nó phải viết lại A, B, C trong khi đầu vào là toàn số 0.}
  \label{fig:ch02-copy-schema}
\end{figure}

Cụ thể, với tham số $T_{\text{mem}}$ (hình~\ref{fig:ch02-copy-schema}):
\begin{enumerate}
  \item $T_{\text{mem}}$ bước đầu: mỗi bước là một vector 8-bit one-hot (8
    lớp, đúng một bit bằng 1).
  \item Một bước \textbf{delimiter}: bit dữ liệu bằng 0, kênh đánh dấu bằng 1.
  \item $T_{\text{mem}}$ bước tiếp theo: toàn bộ đầu vào bằng 0. Mạng phải
    xuất ra đúng $T_{\text{mem}}$ vector đã đọc.
\end{enumerate}

Tổng chiều dài chuỗi: $2T_{\text{mem}} + 1$. Bài toán khó hơn bài toán cộng
vì mạng phải nhớ \textbf{toàn bộ} chuỗi đầu vào qua $T_{\text{mem}}$ bước
hoàn toàn im lặng; không một gợi ý nào từ input.

\subsection{Kết quả}

Tôi huấn luyện RNN với $d_h = 32$, learning rate $0.01$, 2000 epoch, khởi tạo
$\mathcal{N}(0, 0.1)$. Hàm mất mát: cross-entropy trên từng bit đầu ra, chỉ
tính ở $T_{\text{mem}}$ bước cuối. Loss của dự đoán ngẫu nhiên trên 8 lớp là
$-\ln(1/8) \approx 2.08$ mỗi bước.

\begin{measured}{Copy task, $d_h = 32$, 2000 epoch}
\small
\begin{tabular}{@{}lrrrr@{}}
\toprule
$T_{\text{mem}}$ & Loss trước & Loss sau &
$\|\partial L / \partial \mat{W}_{hh}\|$ tại $t=1$ & tại $t=2T_{\text{mem}}+1$ \\
\midrule
5  & 8.74 & 8.34 & 0.233 & 0.033 \\
10 & 7.49 & 8.34 & 0.416 & 0.017 \\
20 & 7.91 & 8.37 & 0.832 & 0.033 \\
\bottomrule
\end{tabular}
\normalsize
\end{measured}

Điểm mấu chốt nằm ở cột cuối: tỉ lệ norm gradient giữa bước đầu tiên và bước
cuối cùng tăng từ 7 lần ở $T_{\text{mem}} = 5$ lên 25 lần ở $T_{\text{mem}}
= 20$. Gradient tại bước xa output yếu hơn gradient tại bước gần output 25
lần; và đây không phải là suy luận, mà là số đo.

Ở $T_{\text{mem}} = 20$, loss sau huấn luyện (8.37) cao hơn loss trước
huấn luyện (7.91). Mạng không những không học được, mà còn tệ hơn cả trước
khi học: gradient yếu đến mức các bước cập nhật trọng số thực chất là nhiễu,
và nhiễu đó phá hỏng cả những gì mạng có từ khởi tạo.

\subsection{Gradient suy giảm theo từng bước}

Bảng trên chỉ so sánh hai điểm: bước đầu và bước cuối. Nhưng gradient không
biến mất đột ngột; nó suy giảm dần qua từng bước lùi. Hình~\ref{fig:ch02-gradient-decay}
cho thấy điều này.

\begin{figure}[ht]
  \centering
  % ch02-gradient-decay-curve.tex — Gradient norm decay curve (schematic)
%
% Stylised log-scale plot showing ||contrib_hh|| decaying exponentially
% as BPTT walks backward from step T to step 1.  Three curves for T_mem=5,10,20.

\begin{tikzpicture}[
    >=Stealth,
    every node/.style={font=\footnotesize},
  ]

  % Axes
  \draw[->] (0,0) -- (8.5,0) node[right] {Bước lùi $T \to 1$};
  \draw[->] (0,0) -- (0,5.5) node[above] {$\|\text{contrib\_hh}\|$ (log)};

  % Y-axis labels (log scale)
  \node[left] at (0,0) {$10^{-4}$};
  \node[left] at (0,1.2) {$10^{-3}$};
  \node[left] at (0,2.4) {$10^{-2}$};
  \node[left] at (0,3.6) {$10^{-1}$};
  \node[left] at (0,4.8) {$10^{0}$};

  % X-axis labels
  \node[below] at (0,0) {$T$};
  \node[below] at (8,0) {$1$};
  \node[below=0.8cm] at (4,0) {$\longleftarrow$ càng xa output, gradient càng yếu};

  % T_mem = 5: shorter, less decay
  \draw[thick, black!70]
    (0,4.5) .. controls (1.5,4.2) and (3.5,3.5) .. (5,2.8)
    .. controls (6,2.2) and (7,1.5) .. (8,1.0);
  \node[right, black!70] at (8,1.0) {$T_{\text{mem}}=5$};

  % T_mem = 10: medium decay
  \draw[thick, black!50, dashed]
    (0,4.5) .. controls (2,4.0) and (4,3.0) .. (6,1.8)
    .. controls (7,1.0) and (7.5,0.5) .. (8,0.15);
  \node[right, black!50] at (8,0.15) {$T_{\text{mem}}=10$};

  % T_mem = 20: strongest decay
  \draw[thick, black!30, dotted]
    (0,4.5) .. controls (2.5,3.5) and (5,1.5) .. (7,0.2)
    .. controls (7.5,0.05) .. (8,0.02);
  \node[right, black!30] at (8,0.02) {$T_{\text{mem}}=20$};

  % Annotation
  \draw[<-] (3.5, 3.8) -- (5, 4.5)
    node[right, font=\scriptsize, text width=3.5cm]
    {Gradient mạnh ở gần output (bên phải)};
  \draw[<-] (6.0, 0.6) -- (7.2, 1.8)
    node[right, font=\scriptsize, text width=3.5cm]
    {Gần như biến mất ở xa output (bên trái)};

\end{tikzpicture}

  \caption{Norm của đóng góp từ mỗi bước vào $\partial L / \partial
    \mat{W}_{hh}$ khi BPTT đi lùi từ $T$ về $1$. Trục tung là thang log.
    Đường cong càng dốc khi $T_{\text{mem}}$ càng lớn: với $T_{\text{mem}} = 20$,
    gradient ở bước $t=1$ nhỏ hơn gradient ở bước $t=T$ khoảng hai bậc độ lớn.}
  \label{fig:ch02-gradient-decay}
\end{figure}

Ba đường cong cho thấy một quy luật: gradient suy giảm theo \textbf{hàm mũ}
của khoảng cách thời gian, không phải tuyến tính. $T_{\text{mem}} = 5$ cho
đường cong thoải; $T_{\text{mem}} = 20$ cho đường cong dốc đứng; gradient
gần như biến mất hoàn toàn sau khoảng 15 bước lùi. Đây chính là lý do RNN
không học được phụ thuộc dài hạn: tín hiệu học từ tương lai không đến được
quá khứ xa.

\index{copy task|)}%
